% II. Анализ на възможните решения и обосновка на избора.
% Тук се прави сравнението „остаряла технология от силабуса“ ⇄ „стандартизиран наследник“.
\section{Анализ на възможните решения и обосновка на избора}
\label{sec:alternatives}

\subsection{Протокол за поток от данни}

Разгледани са четири варианта за доставяне на нови стойности до браузъра: периодично
запитване, задържано запитване, събития, изпращани от сървъра, и
WebSocket~\cite{rfc6455}. Първите два остават в модела „заявка, отговор“ и плащат за
всяка стойност цената на пълен обмен от заглавни части. Третият премахва тази цена, но е
еднопосочен: клиентът няма как да смени наблюдавания набор от показатели по същия канал и
му е нужна отделна заявка за всяко управляващо действие.

Избран е WebSocket, защото задачата изисква и двете посоки: сървърът изпраща стойности,
клиентът изпраща команди за смяна на набора, за смяна на честотата и за смяна на формата
на обмена. Събитията, изпращани от сървъра, остават реализирани като втори, резервен
канал, което позволява двата подхода да бъдат измерени един срещу друг върху един и същ
поток от данни, вместо да бъдат сравнявани по литературни твърдения.

Решението е резервният канал да влезе в предадената версия. Реализирани са и трите
подхода: WebSocket, събития, изпращани от сървъра, и периодично запитване. Клиентът
превключва между тях по време на работа, без презареждане на страницата. Причината е, че
без действаща реализация и на трите сравнението в Раздел~\ref{sec:results} щеше да
съпоставя измерена стойност с цитирана, а измерените стойности се получават само от
работещ код. Цената е приблизително сто реда в сървъра и три малки модула в клиента.

\subsection{Приложна рамка на клиента}

Съвременната практика решава задачата за свързване на данни с приложна рамка, която сама
следи кои части от документа зависят от кои стойности. Разгледани са два пътя: рамка от
познатите днес семейства и клиент без рамка, който работи направо срещу обектния модел на
документа~\cite{dom_ls}.

Избран е клиентът без рамка. Мотивът е учебен: цялото
съдържание на дисциплината, което се отнася до клиентската част (обектният модел,
колекциите, събитията, свързването на данни, структурната графика), при използване на
рамка се превръща в извикване на нейни функции и остава скрито. Без рамка същите теми са
видими в изходния текст на проекта. Цената на това решение е ръчно написан слой за
свързване на данни, чието проектиране е описано в Раздел~\ref{sec:design}.

\subsection{Съответствие между технологиите от учебната програма и избраните}

Темите от учебната програма, които сочат към технологии, изместени от практиката, са
покрити от стандартизираните им наследници. Съответствието е дадено в таблицата по-долу
\tabref{tab:tech-mapping}, а обосновката за всеки ред следва
логиката от Раздел~\ref{sec:literature}: избраната технология покрива същата нужда, но е
част от стандарт и работи във всеки съвременен браузър.

\begin{table}[H]
\caption{Съответствие между технологиите от учебната програма и реализираните в проекта}
\label{tab:tech-mapping}
\centering
\fontsize{11}{13.2}\selectfont
\begin{tabular}{@{}L{4.2cm}L{4.2cm}L{4.6cm}@{}}
\toprule
\textbf{Тема} & \textbf{Реализирано с} & \textbf{Основание} \\
\midrule
Двоични компоненти в страницата (ActiveX) & стандартни браузърни интерфейси & еднаква работа във всеки браузър, без права на потребителя за страницата \\
Генериране на страници на сървъра (ASP) & генериране на HTTP отговор от сървъра на C++20 & същата роля, без обвързаност с една платформа \\
Филтри и преходи като собствено разширение & каскадни стилови преходи и филтри & стандартизирани и хардуерно ускорени \\
Растерна графика за диаграми & структурна графика SVG~\cite{svg2} & всеки елемент е възел в документа и приема събития \\
Свързване на данни чрез компонент & собствен слой за свързване върху обектния модел~\cite{dom_ls} & темата остава видима в изходния текст \\
\bottomrule
\end{tabular}
\end{table}

\subsection{Формат на обмена}

Двата разгледани формата, разширяемият език за разметка~\cite{xml10} и нотацията за
обекти на JavaScript~\cite{rfc8259}, не се изключват взаимно. Проектът реализира и двата
зад един и същ договор за съобщение и позволява формата да се превключва по време на
работа. Това решение превръща избора между тях от предположение в измерване и осигурява
материала за Раздел~\ref{sec:results}.

\subsection{Език и среда на сървъра}

Сървърът е на C++20~\cite{iso_cpp20}. Езикът е избран заради прякото управление на паметта
и предвидимото време за реакция, което е съществено, когато измерваната величина е
закъснението на самата система. Разгледана е и възможността да се използва готов сървър на
приложения, при която обработката на протокола отпада; тя е отхвърлена, защото
обработката на HTTP/1.1~\cite{rfc9112} и на ръкостискането по WebSocket~\cite{rfc6455} е
част от учебното съдържание, а не спомагателна работа.

\subsection{Външни зависимости}

Разгледани са два пътя: сървърът да стъпи върху готова библиотека за мрежов вход и изход,
или да ползва само стандартната библиотека на езика и интерфейса за сокети на
операционната система.

Избран е вторият. Хранилището няма нито една задължителна външна зависимост. Мотивът е
проверимост: чужд човек трябва да може да изпълни \code{cmake -S . -B build} и след това
\code{cmake --build build} на чиста машина, и изграждането да мине от първия път, без
управител на пакети и без изтегляне на нещо от мрежата при конфигуриране. Всичко, което
иначе би дошло от библиотека, е написано в проекта и е малко по обем: инструментът за
изпълнение на тестовете е сто и деветнадесет реда, измерването на времето стъпва на
\code{std::chrono::steady\_clock}, а SHA-1 и Base64 съществуват в изходния текст само
защото ръкостискането по WebSocket~\cite{rfc6455} ги изисква. Свързват се единствено
библиотеките, които операционната система така или иначе доставя.

Цената на това решение е, че някои части са по-прости от промишлените си съответствия:
анализаторът на разширяемия език за разметка~\cite{xml10} покрива подмножеството, което
проектът излъчва, а не целия стандарт. Ограничението е заявено изрично в
Раздел~\ref{sec:implementation}, за да не бъде прието наум.
